\subsection{BÀI TẬP TRẮC NGHIỆM}
\ind{PHẦN I.} \inden{Câu trắc nghiệm nhiều phương án lựa chọn. Mỗi câu hỏi học sinh chỉ chọn một phương án.}\\
\setcounter{ex}{0}
\Opensolutionfile{ans}[ans/1D7-B3-1]
\begin{ex}%[Trần Hoàng Trọng Ân]%[1D5B5]
	Cho hàm số $ y=3x $. Khi đó, $ y'' $ bằng bao nhiêu?
	\choice
	{$3$}
	{$3x$}
	{$\displaystyle\dfrac{3}{2}x^{2}$}
	{\True $0$}
\end{ex}

\begin{ex}%[Trần Hoàng Trọng Ân]%[1D5B5]
	Cho $ f(x)=(x+10)^{6} $. Tính $ f''(2) $.
	\choice
	{$f''(2) = 827440$}
	{\True $f''(2) = 622080$}
	{$ f''(2) = 0$}
	{$f''(2)=413720$}
\end{ex}

\begin{ex}%[Trần Hoàng Trọng Ân]%[1D5B5]
	Cho $ f(x)=\cos 3x $. Tính $ f''(0) $.
	\choice
	{\True$ f''(0)=-9 $}
	{$ f''(0)= 0$}
	{$ f''(0)=9 $}
	{$ f''(0)= 1$}
\end{ex}

\begin{ex}%[Nguyễn Minh Tiến]%[1D5B5]
	Tìm đạo hàm cấp hai của hàm số $y = 5\sin x - 3\cos x$.
	\choice
	{$y'' = 5\cos x + 3\sin x$}
	{$y'' = 5\sin x + 3\cos x$}
	{$y'' = - 5\sin x - 3\cos x$}
	{\True $y'' = - 5\sin x + 3\cos x$}
\end{ex}


\begin{ex}%[Dự án chuyển word 10-11, Chu Đức Minh]%[1D5B5-2]
	Cho hàm số $y=5-\dfrac{3}{x}$. Tính giá trị biểu thức $M=xy''+2y'$.
	\choice
	{\True $M=0$}
	{$M=1$}
	{$M=4$}
	{$M=10$}
	\loigiai{
		\begin{itemize}
			\item Ta có $y'=\dfrac{3}{x^2}\Rightarrow y''=-\dfrac{6}{x^3}$. 
			\item Khi đó $M=x\left(-\dfrac{6}{x^3}\right)+2\cdot \dfrac{3}{x^2}=-\dfrac{6}{x^2}+\dfrac{6}{x^2}=0$.
		\end{itemize}
	}
\end{ex}

\begin{ex}%[Dự án chuyển word 10-11, Chu Đức Minh]%[1D5B5-2]
	Cho hàm số $y= \cos2x$ có đạo hàm là $y'$ và $y''$. Mệnh đề nào sau đây đúng?
	\choice
	{$y+y''=0$}
	{$4y''-y=0$}
	{\True $y''+4y=0$}
	{$y+2y'=0$}
	\loigiai{
		Ta có $y'=-2\sin 2x\Rightarrow y''=-4\cos 2x=-4y\Rightarrow y''+4y=0$.}
\end{ex}

\begin{ex}%[Dự án chuyển word 10-11, Chu Đức Minh]%[1D5B5-2]
	Cho hai hàm số $f(x)=x^4-4x^2+3$ và $g(x)=3+10x-7x^2$. Nghiệm của phương trình $f''(x)+g'(x)=0$ là
	\choice
	{\True $x=1$; $x=\dfrac{1}{6}$}
	{$x=-1$; $x=\dfrac{1}{6}$}
	{$x=-1$; $x=-\dfrac{1}{6}$}
	{$x=1$; $x=-\dfrac{1}{6}$}
	\loigiai{
		\begin{itemize}
			\item Ta có $\heva{& f'(x)=4x^3-8x \\ & g'(x)=-14x+10} \Rightarrow \heva{& f''(x)=12x^2-8 \\ &g'x=-14x+10.}$
			\item Khi đó, phương trình 
			$$f''(x)+g'(x)=0\Leftrightarrow \left(12x^2-8\right)+\left(-14x+10\right)=0 \Leftrightarrow 12x^2-14x+2=0\Leftrightarrow \hoac{& x=1 \\ & x=\dfrac{1}{6}. }$$
		\end{itemize}
	}
\end{ex}


\begin{ex}%[1D5B5]
	Một chất điểm chuyển động thẳng với phương trình $s(t)=t^3-t^2+5t$ ($s$ tính theo $m$, $t$ tính theo $s$). Tính giá trị gia tốc của chất điểm tại thời điểm $t=1s$?
	\choice
	{$a(3)=5m/s^2$}
	{$a(3)=6m/s^2$}
	{$a(3)=0m/s^2$}
	{\True $a(3)=4m/s^2$}
\end{ex}

\begin{ex}%[1D5B2-6]
	Cho chuyển động thẳng xác định bởi phương trình $s(t)=t^3+4t^2$, trong đó $t>0$, $t$ tính bằng giây và $s(t)$ tính bằng mét. Gia tốc của chuyển động tại thời điểm mà vận tốc của chuyển động bằng $11$ m/s là
	\choice
	{$12$ m/s$^2$}
	{$16$ m/s$^2$}
	{\True $14$ m/s$^2$}
	{$18$ m/s$^2$}
	\loigiai{
		\begin{itemize}
			\item Ta có $v(t)=s'(t)=3t^2+8t\Rightarrow a(t)=v'(t)=6t+8$. 
			\item Ta có $v(t)=11\Leftrightarrow 3t^2+8t=11\Leftrightarrow \hoac{
				& t=1 \\ 
				& t=-\dfrac{11}{3} \text{ (loại).}}$
			\item Với $t=1$, ta có $a(1)=14$ m/s$^2$. 
		\end{itemize}
	}
\end{ex}

\begin{ex}%[1D5T3-3]
	Một chất điểm có phương trình chuyển động $s(t)=3 \sin\left(t + \dfrac{\pi}{3} \right) $, trong đó $t>0, t$ tính bằng giây, $s(t)$ tính bằng centimet. Tính gia tốc của chất điểm tại thời điểm $t= \dfrac{\pi}{2}$ (s).
	\choice
	{$-\dfrac{2\sqrt{3}}{3}$}
	{\True $-\dfrac{3\sqrt{3}}{2}$}
	{$-\dfrac{3\sqrt{2}}{2}$}
	{$-\dfrac{2\sqrt{2}}{3}$}
	\loigiai{
		Ta có $a(t)=s''(t)=3\cos\left(t + \dfrac{\pi}{3} \right) $.\\
		$a\left(\dfrac{\pi}{2} \right) =3\cos\left(\dfrac{\pi}{2} + \dfrac{\pi}{3} \right) = 3\cos \left(\dfrac{5\pi}{6}  \right)=-\dfrac{3\sqrt{3}}{2} $  ($\mathrm{m/s^2}$).
	}
\end{ex}
\Closesolutionfile{ans}

\ind{PHẦN II.} \inden{Câu trắc nghiệm đúng sai. Trong mỗi ý a), b), c), d) ở mỗi câu, học sinh chọn đúng hoặc sai.}\\
\Opensolutionfile{ans}[ans/1D7-B3-2]

\begin{ex}%Câu 4.%[1D7V2-8]
	Một vật chuyển động trên đường thẳng được xác định bởi công thức $s(t)=t^3-3t^2+7t-2$, trong đó $t>0$ và tính bằng giây và $s$ là quãng đường chuyển động được của vật trong $t$ giây tính bằng mét. Khi đó
	\choiceTF
	{\True Tốc độ của vật tại thời điểm $t=2$ là $7$ (m/s)}
	{\True Gia tốc của vật tại thời điểm $t=2$ là $6$ (m/s$^2$)}
	{Gia tốc của vật tại thời điểm mà vận tốc của chuyển động bằng $16$ (m/s) là $10$ (m/s$^2$)}
	{\True Thời điểm $t=1$ (s), tại đó vận tốc của chuyển động đạt giá trị nhỏ nhất}
	\loigiai{
		Ta có $s'(t)=3t^2-6t+7$ và $s''(t)=6t-6$.
		\begin{itemchoice}
			\itemch  \textbf{Đúng}.\\
			Vận tốc của vật tại thời điểm $t=2$ là $v(2)=s'(2)=3\cdot 2^2-6\cdot 2+7=7$ (m/s).
			\itemch  \textbf{Đúng}.\\
			Gia tốc của vật tại thời điểm $t=2$ là $a(2)=v'(2)=s''(2)=6\cdot 2-6=6$ (m/s$^2$).
			\itemch  \textbf{Sai}.\\
			Vận tốc của chuyển động bằng $16$ (m/s$^2$) tại thời điểm $t$ nghĩa là
			$$v(t)=s'(t)=16\Leftrightarrow 3t^2-6t+7=16\Leftrightarrow\hoac{&t=3 (\text{thỏa mãn})\\&t=-1 (\text{loại}).}$$
			Gia tốc của vật tại thời điểm $t=3$ là $a(3)=v'(3)=s''(3)=6\cdot 3-6=12$ (m/s$^2$).
			\itemch  \textbf{Đúng}.\\
			Vận tốc của chuyển động có phương trình $v(t)=3t^2-6t+7$ là một parabol, có đỉnh là
			$S\left(-\dfrac{b}{2a};-\dfrac{\Delta}{4a}\right)\Rightarrow S(1;4)$ và hệ số $a=3>0$ nên hàm số có giá trị nhỏ nhất bằng 4 tại $t=1$
			(hoặc sử dụng máy tính Casio để tìm min).\\
			Vậy tại thời điểm $t=1$ thì vận tốc của chuyển động đạt giá trị nhỏ nhất bằng $4$ (m/s).
	\end{itemchoice}}
\end{ex}

\begin{ex}
	Cho hàm số $y=\dfrac{x+2}{x-1}(C)$.
	\choiceTF
	{\True $y'(0)=-3$}
	{\True $y''=\dfrac{6}{(x-2)^3}$}
	{Tiếp tuyến tại điểm có hoành độ $x=0$ có phương trình $y=-3x+2$}
	{\True Nếu $a \leq-2$ thì qua điểm $A(0 ; a), a \in \mathbb{R}$ sẽ kẻ được hai tiếp tuyến tới $(C)$}
	\loigiai{
		\begin{itemchoice}
			\itemch Đúng: Ta có $y'=\dfrac{-3}{(x-1)^2} \Rightarrow y'(0)=-3$.
			\itemch  Đúng: Ta có $y'=\dfrac{-3}{(x-1)^2} \Rightarrow y''=\dfrac{3 \cdot 2 \cdot(x-1)'(x-1)}{(x-1)^4}=\dfrac{6}{(x-1)^3}$.
			\itemch Sai: Với $x=0$ thì $y(0)=-2$ và $y'(0)=-3$.\\
			Tiếp tuyến tại điểm có hoành độ $x=0$ có phương trình $y=-3(x-0)-2 \Rightarrow y=-3 x-2$.
			\itemch Sai: Ta có: $y'=f'(x)=\dfrac{-3}{(x-1)^2}$.\\
			Phương trình tiếp tuyến d qua $A(0 ; a), a \in \mathbb{R}:(d): y=\dfrac{-3}{(x-1)^2} x+a$.\\
			Lại có $d$ tiếp xúc $(C)$ nên xét phương trình hoành độ giao điểm giữa $d$ và $(C)$
			$$
			\frac{x+2}{x-1}=\frac{-3}{(x-1)^2} x+a \Leftrightarrow
			\heva{&x \neq 1 \\
				&	(1-a) x^2+2(a+2) x-(a+2)=0.}
			$$
			Ta đặt vế trái bằng $f(x)$. Qua $A$ kẻ được 2 tiếp tuyến thì 
			$f(x)=0$ có hai nghiệm phân biệt $x_1$; $x_2$ và khác $1$\\
			$
			\Leftrightarrow\heva{&1 - a \neq 0 \\&\Delta'= (a+2)^2+(1-a)(a+2)>0\\&f(1)=(1-a) \cdot 1^{ 2}+2(a+2) \cdot 1-(a+2) \neq 0} \Leftrightarrow \heva{&a \neq 1\\&a > - 2\\&3 \neq 0} \Leftrightarrow \heva{&a \neq 1 \\&a>-2.}
			$
		\end{itemchoice}
	}
\end{ex}

\Closesolutionfile{ans}

\ind{PHẦN III.} \inden{Câu trắc nghiệm trả lời ngắn.}\\
\Opensolutionfile{ans}[ans/1D7-B3-3]

\begin{ex}%[1D7N3-1] % Câu 9:
	Cho hàm số $f(x)=\mathrm{e}^{2x}$. TÍnh $f''(1)$ (kết quả làm tròn đến hàng phần chục). 
	\\ \shortans[oly]{$29,6$}
	\loigiai{
		Ta có $f(x)=\mathrm{e}^{2x} \Rightarrow f'(x)=2\mathrm{e}^{2x} \Rightarrow f''(x)=4\mathrm{e}^{2x}$.\\
		Vậy $f''(1)=4\mathrm{e}^2 \approx 29,6$.}
\end{ex}

\begin{ex}%[1D7N3-1] % Câu 10
	Cho hàm số $f(x)=\log_2(3x+1)$. Tính $f''(2)$ (kết quả làm tròn đến hàng phần chục). 
	\\ \shortans[oly]{$-0,3$}
	\loigiai{
		Ta có $f(x)=\log_2(3x+1) \Rightarrow f'(x)=\dfrac{3}{(3x+1)\ln 2}$.\\
		$ \Rightarrow f''(x)=\dfrac{-9}{(3x+1)^2\ln 2}$.\\
		Vậy $f''(2)=\dfrac{-9}{49\ln 2} \approx -0,3$.}
\end{ex}

\begin{ex}
	Một vật chuyển động thẳng không đều xác định bởi phương trình $s(t)=2 t^{2} + 5 t + 9$, trong đó ${s}$ tính bằng mét và ${t}$ tính bằng giây. Tính gia tốc của chuyển động tại thời điểm $t=2$.\\
	\\ \shortans[oly]{$4$}
	\loigiai{ 
		
		$v(t)=\left(2 t^{2} + 5 t + 9\right)'=4 t + 5.$
		
		$a(t)=\left(4 t + 5\right)'=4$.
		
		$a(2)=4$ $m/s^2$. 
}\end{ex}

\begin{ex}%Câu 5.%[1D7V2-8]
	Một vật chuyển động thẳng được xác định bởi phương trình $s(t)=t^3-2t^2+2t$, trong đó $t$ tính bằng giây và $s$ tính bằng mét. Tính gia tốc của vật tại thời điểm mà vận tốc của vật bằng $17$ (m/s).\\
	\shortans[oly]{$14$}
	\loigiai{
		Ta có vận tốc của vật tại thời điểm $t$ là $v(t)=s'(t)=3t^2-4t+2$.\\
		Gia tốc của vật tại thời điểm $t$ là $a(t)=s''(t)=6t-4$.\\
		Vận tốc của vật bằng $17$ (m/s) tại thời điểm $t$ suy ra $$v(t)=17\Leftrightarrow 3t^2-4t+2=17\Leftrightarrow\hoac{&t=3 (\text{thỏa mãn})\\&t=-\dfrac{5}{3}\quad(\text{loại}).}$$
		Gia tốc của vật tại thời điểm $t=3$ là $a(3)=6\cdot 3-4=14$ (m/s$^2$).}
\end{ex}

\Closesolutionfile{ans}

